\section{Finite Hurwitz identities and valid amplification}
\label{sec:hurwitz-amplification}

A short identity at one dimension need not give an asymptotic saving.
This issue matters for the complex formula of type $[12,12,18]$ of
Zhang--Zhu \cite[Theorem~1.1, equations~(3.1)--(3.3)]{ZhangZhu2026}:
\[
 \Bigl(\sum_{i=1}^{12}x_i^2\Bigr)
 \Bigl(\sum_{j=1}^{12}y_j^2\Bigr)
   =\sum_{\alpha=1}^{18}z_\alpha(x,y)^2.
\]
The squares are algebraic squares, without complex conjugation.
The displayed coefficient array in that source passes all $12^4=20736$
polarized coefficient equations in our exact Gaussian-integer check.
The following results analyze two valid ways to amplify it. Neither
method reaches the exponent $\sqrt2$ relevant to the square.

\subsection{Tensor products with a parity correction}

Write a formula of this type as Hurwitz--Radon matrices
$H_1,\ldots,H_{12}\in\mathbb C^{12\times18}$ satisfying
\[
 H_iH_i^{\mathsf T}=I_{12},\qquad
 H_iH_j^{\mathsf T}+H_jH_i^{\mathsf T}=0\quad(i\ne j).
\]
For a word $a\in[12]^\ell$ put
$H_a=H_{a_1}\otimes\cdots\otimes H_{a_\ell}$.
These matrices do not themselves satisfy the Hurwitz--Radon conditions:
a product of an even number of nonzero skew matrices, on separate
factors, is symmetric. A valid construction adds a vector $\xi_a$ to
correct this parity; see \cite[Section~5.2, p.~12]{Hrubes2026}.

\begin{proposition}[Cost of the full-word parity correction]
\label{prop:hurwitz-parity-cost}
For each $a\in[12]^\ell$, choose $\xi_a\in\mathbb C^t$ with
$\xi_a^{\mathsf T}\xi_a=1$ and put
$K_a=H_a\otimes\xi_a^{\mathsf T}$.
The matrices $K_a$ define a Hurwitz formula if and only if
\[
 \xi_a^{\mathsf T}\xi_b=0
 \quad\hbox{whenever }a\ne b\hbox{ have even Hamming distance}.
\]
Consequently
\[
 t\ge12^{\lfloor\ell/2\rfloor}.
\]
Its declared output length $18^\ell t$ has asymptotic exponent at
least $\log_{12}18+1/2$. Moreover the Gram matrix of its coefficient
vectors has rank at least
$12^{\ell+\lfloor\ell/2\rfloor}$; a compression preserving that Gram
matrix therefore still has exponent at least $3/2>\sqrt2$.
\end{proposition}

\begin{proof}
Each $H_i$ has row rank twelve. Sylvester's rank inequality gives
$\operatorname{rank}(H_iH_j^{\mathsf T})\ge12+12-18=6$,
so all these products are nonzero. Thus
$H_aH_b^{\mathsf T}$ is nonzero, and its transpose is itself times
$(-1)^{d_H(a,b)}$. Since
\[
 K_aK_b^{\mathsf T}
   =(H_aH_b^{\mathsf T})(\xi_a^{\mathsf T}\xi_b),
\]
the off-diagonal Hurwitz--Radon condition is automatic at odd
Hamming distance and requires exactly the stated orthogonality at
even distance.

Consider words whose consecutive coordinate pairs are equal, fixing
the last coordinate when $\ell$ is odd. There are
$12^{\lfloor\ell/2\rfloor}$ such words, and any two distinct ones have
positive even Hamming distance. Their vectors $\xi_a$ are mutually
orthogonal of norm squared one, hence linearly independent. This
proves the bound on $t$.

For each word in this set, all $12^\ell$ rows of $K_a$ are
orthonormal. Rows from different words in the set are orthogonal
because their correction vectors are. These rows form an identity
principal submatrix of the coefficient-vector Gram matrix of size
$12^{\ell+\lfloor\ell/2\rfloor}$. This proves the rank bound and its
consequence for Gram-preserving compression.
\end{proof}

Hrube\v{s}'s construction supplies correction vectors with
$t\le48^{\ell/2}$ \cite[Section~5.2, p.~12]{Hrubes2026}. Thus its
stated upper exponent for this seed is
\[
 \log_{12}(18\sqrt{48})=1.9421141086\ldots,
\]
whereas the declared-length lower bound in the proposition is
\[
 \log_{12}18+1/2=1.6631711630\ldots.
\]
The proposition concerns all tensor words with this scalar-vector
correction. It does not constrain a different composition rule,
a selected subfamily of words, or a recompression that changes the
Gram matrix while preserving only the polynomial identity.

The failure of uncorrected tensor powering is already explicit in
the Zhang--Zhu table. With one-based indices, put
$R=H_1H_2^{\mathsf T}$. Its entries satisfy $R_{12}=1$ and
$R_{21}=-1$. For the word pair $(1,1),(2,2)$, the corresponding
Hurwitz anticommutator has entry
$R_{12}^2+R_{21}^2=2$ at row pair $(1,1),(2,2)$, where zero is
required. This disproves that particular tensor-power construction,
not the existence of other formulas at the enlarged dimensions.

\subsection{Majority composition under partial transpose}

For a $q^2$ by $q^2$ matrix, write
\[
 (\Gamma A)_{ij,k l}=A_{i l,kj}
\]
for partial transpose, with $i,j,k,l\in[q]$. If $V$ is the coefficient
matrix of a complex bilinear sum-of-squares identity, polarization
gives $A+\Gamma A=I_{q^2}$ for $A=VV^{\mathsf T}/2$.
This admits a composition with a directly calculable rank, without
asserting submultiplicativity of general Hurwitz length.

\begin{lemma}[Rank of a complementary partial-transpose seed]
\label{lem:partial-transpose-seed-rank}
Over a field of characteristic different from two, every matrix
$A$ satisfying $A+\Gamma A=I_{q^2}$ has rank at least $q$.
No symmetry assumption on $A$ is needed.
In characteristic two this identity has no solution.
\end{lemma}

\begin{proof}
Fix a second-factor index $j$. Partial transpose fixes every entry
of the principal block with rows and columns $(i,j)$, so
\[
 2A_{(i,j),(k,j)}=\delta_{ik}.
\]
That block is $I_q/2$, proving the rank bound. In characteristic
two the diagonal entries instead give the contradiction $0=1$.
\end{proof}

\begin{theorem}[Exact rate of majority composition]
\label{thm:hurwitz-majority-rate}
Let $q\ge2$, put $d=q^2$, and suppose
$A+\Gamma A=I_d$ and
$\operatorname{rank}A=\operatorname{rank}A^2=r<d/2$ over
$\mathbb C$. For odd $k$ let $B_k$ be the sum of tensor words in
$A,I-A$ containing strictly more than $k/2$ factors $A$.
Then $B_k+\Gamma B_k=I_{d^k}$, where partial transpose respects the
grouping into two $q^k$-dimensional factors. With $j_0=(k+1)/2$,
\begin{equation}\label{eq:majority-rank-bounds}
 \binom{k}{j_0}r^{j_0}(d-r)^{k-j_0}
 \le\operatorname{rank}B_k
 \le\sum_{j=j_0}^k\binom{k}{j}r^j(d-r)^{k-j}.
\end{equation}
In particular
\[
 \lim_{\substack{k\to\infty\\k\ \text{odd}}}
       (\operatorname{rank}B_k)^{1/k}=2\sqrt{r(d-r)}.
\]
If $A$ is symmetric, $B_k$ supplies a complex sum-of-squares formula
of type $[q^k,q^k,\operatorname{rank}B_k]$.
For every seed in the theorem,
\[
 2\sqrt{r(q^2-r)}\ge2q\sqrt{q-1}>q^{3/2}>q^{\sqrt2}.
\]
Thus this family has exponent strictly above the square threshold,
even when the fixed seed is optimized under the stated hypotheses.
\end{theorem}

\begin{proof}
Partial transpose exchanges $A$ and $I-A$ and acts separately on
the tensor factors. For odd $k$ it pairs every majority word with
one minority word. Their total is
$(A+(I-A))^{\otimes k}=I_{d^k}$, proving the identity for $B_k$.

The rank hypothesis gives a direct sum
$\mathbb C^d=U\oplus W$, where $U=\operatorname{im}A$,
$W=\ker A$, $\dim U=r$ and $\dim W=d-r$.
Both $A$ and $I-A$ preserve this decomposition; $A|_U$ is invertible,
and $I-A$ is the identity on $W$. Decompose the tensor power by the
set of coordinates belonging to $U$. On a block with fewer than
$j_0$ such coordinates every majority word vanishes. On a block
with exactly $j_0$ such coordinates, precisely one majority word
survives, and it is invertible there. The other blocks have at most
their full dimension as rank. Summing these dimensions gives
\eqref{eq:majority-rank-bounds}.

The ratio of successive summands of the upper bound, starting at
$j_0$, is at most $r/(d-r)<1$. Hence the upper bound is at most its
first term times $1/(1-r/(d-r))$. Since the central binomial
coefficient has $k$th root tending to two, the two bounds have the
stated common rate.

If $A$ is symmetric, so is $B_k$. A complex symmetric matrix of rank
$R$ has a factorization as a sum of $R$ outer squares. For a product
vector $x\otimes y$, partial transpose does not change its quadratic
value. The identity for $B_k$ therefore makes the quadratic form
$2B_k$ equal to $(x^{\mathsf T}x)(y^{\mathsf T}y)$ on product vectors,
which is the required bilinear sum-of-squares identity.

Finally Lemma~\ref{lem:partial-transpose-seed-rank} gives $r\ge q$.
The function $r(d-r)$ is increasing for $r<d/2$, proving the
asserted lower bound.
\end{proof}

For the Zhang--Zhu formula the first three output coordinates give
\[
 \begin{pmatrix}12&0&-8i\\0&8&0\\-8i&0&-8\end{pmatrix}
\]
as their block of $V^{\mathsf T}V$. The other coordinates give
$12I_3\oplus8I_{12}$, with zero off-block entries. The first block
has eigenvalues $8,8,-4$. Consequently $A=VV^{\mathsf T}/2$ has
nonzero eigenvalues $6$ with multiplicity three, $4$ with multiplicity
fourteen, and $-2$ with multiplicity one. In particular
$\operatorname{rank}A=\operatorname{rank}A^2=18$.
The exact majority exponent is therefore
\[
 \log_{12}\bigl(2\sqrt{18\cdot126}\bigr)
       =1.8336600344\ldots>\sqrt2.
\]
The theorem permits arbitrary cancellations within these majority
sums: its lower bound comes from blocks on which precisely one word
survives. Its exact rate assumes rank stability at zero. The following
lower bound removes that hypothesis, while retaining the restriction
that the seed rank is less than half the ambient dimension.

\begin{lemma}[A compression bound for majority sums]
\label{lem:majority-compression}
Let $A\in\mathbb C^{d\times d}$ have rank $r$, and define $B_k$ as
above for odd $k$. If $1\le m\le r$ and $w=d-r-m>0$, then
\[
 \operatorname{rank}B_k\ge
 \binom{k}{(k+1)/2}m^{(k+1)/2}w^{(k-1)/2}.
\]
In particular
\[
 \liminf_{\substack{k\to\infty\\k\ \text{odd}}}
       (\operatorname{rank}B_k)^{1/k}\ge2\sqrt{m(d-r-m)}.
\]
No partial-transpose identity or rank-stability assumption is needed
for this bound.
\end{lemma}

\begin{proof}
Choose an $m$-dimensional subspace $U$ on which $A$ is injective,
and put $J=A(U)$. The restriction of $A$ to $U+J$ has rank at least
$m$, so
\[
 \dim((U+J)\cap\ker A)\le\dim(U+J)-m\le m.
\]
Since $\dim\ker A=d-r$, there is a $w$-dimensional subspace
$W\subseteq\ker A$ with $W\cap(U+J)=0$.
Identify $U$ and $J$ through $A|_U$, and let
$R:\mathbb C^m\oplus\mathbb C^w\longrightarrow\mathbb C^d$
include $U\oplus W$.
Choose a linear map
$L:\mathbb C^d\longrightarrow\mathbb C^m\oplus\mathbb C^w$
whose first component is the chosen identification on $J$ and
vanishes on $W$, and whose second component is the identification
on $W$ and vanishes on $U+J$. These choices are compatible because
the indicated subspaces are disjoint. Then, for some $m$ by $m$
matrix $M$,
\[
 LAR=\begin{pmatrix}I_m&0\\0&0\end{pmatrix},\qquad
 L(I-A)R=\begin{pmatrix}M&0\\0&I_w\end{pmatrix}.
\]
Apply $L^{\otimes k}$ on the left and $R^{\otimes k}$ on the
right of $B_k$. The resulting operator is block diagonal according
to which coordinates use the $m$-dimensional summand. On every
block with exactly $(k+1)/2$ such coordinates, precisely one
majority word survives, and it is the identity on that block.
These blocks have total dimension equal to the asserted lower
bound. Compression cannot increase rank. Taking $k$th roots proves
the limiting inequality.
\end{proof}

\begin{corollary}[The majority barrier with arbitrary zero blocks]
\label{cor:majority-jordan-barrier}
Suppose $q\ge2$, $d=q^2$, $A+\Gamma A=I_d$, and
$r=\operatorname{rank}A<d/2$. With
\[
 m=\min\{r,\lfloor(d-r)/2\rfloor\},\qquad w=d-r-m,
\]
one has
\[
 \liminf_{\substack{k\to\infty\\k\ \text{odd}}}
       (\operatorname{rank}B_k)^{1/k}
       \ge2\sqrt{mw}>q^{\sqrt2}.
\]
For $q\ge4$, the displayed lower bound is even strictly larger than
$q^{3/2}$. Thus a fixed seed of rank less than $q^2/2$ cannot yield
a majority exponent below $\sqrt2$, including when its zero
eigenvalue has nontrivial Jordan blocks.
\end{corollary}

\begin{proof}
Lemma~\ref{lem:partial-transpose-seed-rank} gives $r\ge q$ without
any rank-stability assumption. The displayed choice of $m$ is
admissible in Lemma~\ref{lem:majority-compression}.

First suppose $q\ge4$ and $r\le d/3$. Then $m=r$ and
$mw=r(d-2r)$. This concave quadratic on $[q,d/3]$ is at least
the smaller of its endpoint values. Hence
\[
 2\sqrt{mw}\ge
 \min\{2q\sqrt{q-2},\tfrac23q^2\}>q^{3/2}.
\]
If instead $r>d/3$, put $s=d-r>d/2$. Now
$m=\lfloor s/2\rfloor$, $w=\lceil s/2\rceil$, and
\[
 2\sqrt{mw}\ge\sqrt{s^2-1}>d/2\ge q^{3/2}.
\]
The strict middle inequality follows from the integrality of $s$
and $d$. For $q=2$ the hypotheses $r\ge q$ and $r<q^2/2$ are
incompatible. For $q=3$ the possible ranks $r=3,4$ give respectively
$6,\sqrt{24}$. Both exceed
$3^{\sqrt2}$: indeed $\sqrt2<10/7$, while
$24^7>3^{20}$ follows from $2^{21}>3^{13}$.
\end{proof}

These bounds concern the specified majority sums. They do not cover
seeds of rank at least half the ambient dimension, a different
amplification rule, or replacement of the coefficient matrix by a
different one with the same polynomial restriction to product vectors.

\begin{proposition}[Complementary similar seeds]
\label{prop:majority-similar-seed}
If $A\in\mathbb C^{d\times d}$ is similar to $I-A$, then every odd
majority sum satisfies
\[
 \operatorname{rank}B_k\ge d^k/2,
 \qquad
 \lim_{\substack{k\to\infty\\k\ \text{odd}}}
       (\operatorname{rank}B_k)^{1/k}=d.
\]
In particular every seed on $\mathbb C^q\otimes\mathbb C^q$ with
$q=2$ and $A+\Gamma A=I_4$ has majority rate four, with no rank or
Jordan-structure restriction. More generally, the same conclusion
with $d=2a$ holds on $\mathbb C^a\otimes\mathbb C^2$ when partial
transpose is taken on the two-dimensional factor.
\end{proposition}

\begin{proof}
If $SAS^{-1}=I-A$, conjugation by $S^{\otimes k}$ exchanges the
majority and minority sums. Their sum is $I_{d^k}$ for odd $k$.
Thus $B_k$ and $I-B_k$ have equal rank, and rank subadditivity gives
$d^k\le2\operatorname{rank}B_k$. The trivial upper bound $d^k$
then proves the limit.

For the rectangular assertion, write $A=I_{2a}/2+K$.
Each two-by-two block of $K$ is skew-symmetric, because
$\Gamma K=-K$. Therefore
\[
 K=M\otimes J,\qquad
 J=\begin{pmatrix}0&1\\-1&0\end{pmatrix}
\]
for an $a$ by $a$ matrix $M$. The matrix
$D=\operatorname{diag}(1,-1)$ satisfies $DJD^{-1}=-J$.
Consequently $I_a\otimes D$ conjugates $A$ to $I-A$, and the
first part applies.
\end{proof}

\subsection{Multilinear recombinations of one seed}

Majority composition is one choice of coefficients in the tensor
words of $A$ and $I-A$. The next bound permits arbitrary complex
coefficients, which may change with the number of copies.

\begin{proposition}[A barrier for all multilinear recombinations]
\label{prop:hurwitz-multilinear-barrier}
Suppose $A+\Gamma A=I_d$, where $d=q^2$, and
$\operatorname{rank}A=\operatorname{rank}A^2=r$ with $0<r<d/2$.
Put $w=d-r$. For each $k\ge1$ let
\[
 B=\sum_{S\subseteq[k]}c_S
       \bigotimes_{i=1}^k Z_i(S),\qquad
 Z_i(S)=\begin{cases}A&i\in S,\\I-A&i\notin S,\end{cases}
 \qquad c_S+c_{[k]\setminus S}=1.
\]
Then
\[
 \operatorname{rank}B\ge r^{\lceil k/2\rceil}w^{\lfloor k/2\rfloor}.
\]
The coefficient condition describes precisely all matrices in
$\operatorname{span}\{A,I-A\}^{\otimes k}$ satisfying
$B+\Gamma B=I_{d^k}$. Thus every sequence in this class has lower
exponential rank rate at least $\sqrt{rw}$.

For the Zhang--Zhu seed this lower rate is
\[
 \sqrt{18\cdot126}=18\sqrt7>12^{3/2}>12^{\sqrt2}.
\]
In particular arbitrary multilinear recombination of this fixed seed
cannot reach the square exponent.
\end{proposition}

\begin{proof}
Since $A$ is neither zero nor an invertible scalar matrix, $A$ and
$I-A$ are linearly independent. Their tensor words are therefore
linearly independent. Partial transpose interchanges the word indexed
by $S$ with that indexed by its complement, while $I_{d^k}$ is the sum
of all the words. This proves the asserted characterization.

Choose a set $S$ of minimum cardinality with $c_S\ne0$, and put
$t=|S|$. A pair of complementary subsets of sizes
$\lfloor k/2\rfloor,\lceil k/2\rceil$ has coefficients summing to
one. Consequently $t\le\lceil k/2\rceil$.

As in Theorem~\ref{thm:hurwitz-majority-rate}, decompose
$\mathbb C^d=U\oplus W$ with $U=\operatorname{im}A$ and
$W=\ker A$. Consider the tensor block with $U$ in exactly the
coordinates in $S$ and $W$ elsewhere. A word with an $A$ factor
outside $S$ vanishes there. A word whose set is properly contained
in $S$ has zero coefficient by minimality. Thus precisely the word
indexed by $S$ survives, acting invertibly on this block with the
nonzero coefficient $c_S$. Its dimension is $r^t w^{k-t}$.
Because $r<w$ and $t\le\lceil k/2\rceil$, this is at least the
stated bound. The numerical seed comparison follows by squaring:
$18^2\cdot7=2268>1728=12^3$.
\end{proof}

\begin{corollary}[Optimality among coefficients invariant under copy permutations]
\label{cor:hurwitz-symmetric-optimal}
Under the hypotheses of Proposition~\ref{prop:hurwitz-multilinear-barrier},
suppose $k$ is odd and $c_S$ depends only on $|S|$. Then, with
$j_0=(k+1)/2$,
\[
 \operatorname{rank}B\ge
 \min\left\{w^k,\binom{k}{j_0}r^{j_0}w^{k-j_0}\right\}.
\]
If $r\le d/5$, the minimum possible rank in this class, with
coefficients optimized separately at each odd $k$, has exponential
rate exactly $2\sqrt{rw}$. Majority attains this rate. In particular
the majority rate computed above is optimal in this class for the
Zhang--Zhu seed.
\end{corollary}

\begin{proof}
Let $t$ be the smallest index with nonzero coefficient. Every set of
size $t$ now supplies an invertible block in the preceding proof, so
\[
 \operatorname{rank}B\ge\binom{k}{t}r^t w^{k-t},\qquad0\le t\le j_0.
\]
The sequence $\binom{k}{j}r^j w^{k-j}$ is log-concave in $j$;
its minimum on the integer interval $[0,j_0]$ is therefore at an
endpoint. This gives the finite bound. Its lower exponential rate is
$\min\{w,2\sqrt{rw}\}$. If $r\le d/5$, then $w\ge4r$, so this
minimum is $2\sqrt{rw}$. The matching upper rate follows from the
majority construction in Theorem~\ref{thm:hurwitz-majority-rate}.
\end{proof}

These results concern the multilinear span of a fixed seed. They do
not bound a construction outside that span, a seed changing with the
number of copies, or a replacement of the final coefficient matrix
by another matrix with the same polynomial restriction to product
vectors.

\subsection{Asymmetric decompositions in the flip--cup algebra}
\label{subsec:flip-cup}

The elementary asymmetric decomposition supplies a different possible
starting point. On $H=\mathbb C^q\otimes\mathbb C^q$ let
\[
 F(e_i\otimes e_j)=e_j\otimes e_i,\qquad
 v=\sum_{i=1}^q e_i\otimes e_i,\qquad \Omega=vv^{\mathsf T}.
\]
Then $\Gamma F=\Omega$, $\Gamma\Omega=F$, and
\[
 I=(I-F)+\Gamma\Omega,\qquad
 \operatorname{rank}(I-F)=b:=\frac{q(q-1)}2,\qquad
 \operatorname{rank}\Omega=1.
\]
For $q=4$ this has total cost $7<4^{\sqrt2}$. The following finite
obstruction explains why multiplying this cost does not even give a
valid two-copy decomposition within the full algebra generated by
these three operators. Put
$\mathcal B_q=\operatorname{span}_{\mathbb C}\{I,F,\Omega\}$;
on its tensor powers, $\Gamma$ acts separately in every copy.

\begin{proposition}[Exact self-complementary rank in two copies]
\label{prop:flip-cup-two-copy}
For $q\ge4$, the minimum rank of
$A\in\mathcal B_q\otimes\mathcal B_q$ satisfying
$A+\Gamma A=I_{q^4}$ is
\[
 R=q^2+b^2.
\]
Consequently every asymmetric decomposition
$I_{q^4}=M+\Gamma N$ with
$M,N\in\mathcal B_q\otimes\mathcal B_q$ has
\[
 \operatorname{rank}M+\operatorname{rank}N
 \ge q^2+b^2=(b+1)^2+q-1.
\]
This last inequality is a lower bound on the asymmetric cost; the
first equality computes the self-complementary minimum.
\end{proposition}

\begin{proof}
Write $K=F-\Omega$ and $J=F+\Omega$. These are respectively odd and
even under $\Gamma$. Every self-complementary matrix therefore has
\[
 2A=I\otimes I+xK\otimes I+yI\otimes K
                  +zK\otimes J+wJ\otimes K.
\]
The three simultaneous eigenspaces are the cup line, of dimension
$1$ and eigenvalues $(K,J)=(1-q,1+q)$; the symmetric traceless
space, of dimension $a=q(q+1)/2-1$ and eigenvalues $(1,1)$; and the
antisymmetric space, of dimension $b$ and eigenvalues $(-1,-1)$.
Use symbols $0,+,-$ for these spaces.
The eigenvalues of $2A$ on $++,+-,-+,--$ are, respectively,
\[
 1+x+y+z+w,\quad 1+x-y-z-w,\quad
 1-x+y-z-w,\quad 1-x-y+z+w.
\]
Their block dimensions are $a^2,ab,ab,b^2$. Since
\[
 ab-R=\frac{q(q^2-4q+1)}2>0\qquad(q\ge4),
\]
any matrix of rank at most $R$ must vanish on the first three
blocks. This forces $x=y=-1$ and $z+w=1$, while the last eigenvalue
is $4$. The five remaining eigenvalues, in the order
$00,0+,0-,+0,-0$, are
\[
 -q(q-2),\quad 2q(1-z),\quad2qz,\quad2qz,\quad2q(1-z).
\]
Their nonzero block dimensions sum to $1+a+b=q^2$ when $z=0$ or
$z=1$, and to $1+2(a+b)$ otherwise. This both proves the lower bound
and gives a rank-$R$ matrix. Finally,
$A=(M+N)/2$ is self-complementary for every asymmetric decomposition,
and its rank is at most $\operatorname{rank}M+\operatorname{rank}N$.
\end{proof}


\begin{proposition}[Exact asymmetric cost in two copies]
\label{prop:flip-cup-asymmetric-two-copy}
For every $q\ge4$, the exact minimum of
$\operatorname{rank}M+\operatorname{rank}N$ over
$I_{q^4}=M+\Gamma N$ with
$M,N\in\mathcal B_q\otimes\mathcal B_q$ is
\[
 C=b^2+2q^2-1.
\]
In particular the minimum is $67$ when $q=4$.
\end{proposition}

\begin{proof}
The functional averaging the $+$ and $-$ eigenvalues is invariant
under $\Gamma$, since it takes values $1,0,0$ on $I,F,\Omega$.
Consequently the sum of the four $\{+,-\}^2$ eigenvalues across $M$
and $N$ is $4$. At least one of these blocks survives. Each costs
at least $b^2$, and $2b^2>C$ for $q\ge4$. Thus any pair of cost at
most $C$ has exactly one such block, with eigenvalue $4$.
Its type cannot be $++$, since $a^2>C$.
Exchanging the pair and the two copies reduces the possibilities to
one block of $M$, of type $--$ or $+-$.

Let $(s,t)$ be that pair of signs. The space of operators vanishing
on all four $\{+,-\}^2$ blocks is
\[
 S_0=\operatorname{span}\{
 \Omega\otimes I,\Omega\otimes F,\Omega\otimes\Omega,
 I\otimes\Omega,F\otimes\Omega\}.
\]
The tensor basis $I,F,\Omega$ gives
$S_0\cap\Gamma S_0=\operatorname{span}\{
\Omega\otimes F,F\otimes\Omega\}$.
Every pair in the given case is therefore
\begin{align*}
 M&=(I+sF)\otimes(I+tF)
       +\alpha\Omega\otimes F+\beta F\otimes\Omega,\\
 N&=-s\Omega\otimes I-tI\otimes\Omega-st\Omega\otimes\Omega
       -\alpha F\otimes\Omega-\beta\Omega\otimes F.
\end{align*}
For $s=t=-1$, the four axis blocks associated with each parameter
$z=\alpha$ or $\beta$ have total rank contribution
\[
 (a+b)\mathbf1_{z\ne0}
       +a\mathbf1_{z\ne1}+b\mathbf1_{z\ne-1}\ge a+b.
\]
The two scalar-scalar eigenvalues sum to $q(2-q)\ne0$, so these
blocks together contribute at least one more. Adding the $--$
block proves the lower bound $b^2+2(a+b)+1=C$.
For $s=1,t=-1$, the four axis blocks associated with $\beta$ alone
contribute
\[
 (a+b)\mathbf1_{\beta\ne0}
       +a\mathbf1_{\beta\ne-1}+b\mathbf1_{\beta\ne1}\ge a+b.
\]
Together with the $+-$ block this gives $ab+a+b>C$, by the same
positive difference $q(q^2-4q+1)/2$ used above.
Finally the pair
\[
 M=(I-F)\otimes(I-F),\qquad
 N=\Omega\otimes I+I\otimes\Omega-\Omega\otimes\Omega
\]
attains $C$: its ranks are $b^2$ and $2q^2-1$.
\end{proof}

\begin{proposition}[An all-copy bound in the same algebra]
\label{prop:flip-cup-all-copy}
For all $q\ge2$ and $k\ge1$, every self-complementary
$A\in\mathcal B_q^{\otimes k}$ satisfies
\[
 \operatorname{rank}A\ge b^k.
\]
The same lower bound holds for the total rank of any asymmetric
$I=M+\Gamma N$ with $M,N\in\mathcal B_q^{\otimes k}$.
For each fixed $q\ge5$, this rate is strictly greater than
$q^{\sqrt2}$. Thus arbitrary recombinations in this full tensor
algebra cannot give the desired exponent in those dimensions.
This bound alone leaves the dimensions $q=2,3,4$ open.
Proposition~\ref{prop:full-flip-cup-four} supplies a stronger
all-copy bound that resolves the $q=4$ case.
\end{proposition}

\begin{proof}
Every self-complementary matrix is $I/2$ plus a linear combination
of words in $I,J,K$ with an odd number of $K$ factors. On the
$2^k$ blocks using only $+$ and $-$, both $J$ and $K$ have the
same sign eigenvalue in each coordinate. Every odd word becomes
a nonconstant sign monomial, whose average over these blocks is
zero. The average eigenvalue of $A$ is consequently $1/2$.
At least one such block survives, and each has dimension at least
$b^k$. Averaging an asymmetric pair proves its bound as before.
For $q\ge6$, $(q-1)^2>4q$ gives $b>q^{3/2}>q^{\sqrt2}$.
For $q=5$, the exact inequalities
$10^7>5^{10}$ and $10/7>\sqrt2$ give the same conclusion.
\end{proof}

These bounds impose no condition on decompositions leaving
$\mathcal B_q^{\otimes k}$ and supply no assertion that the optimal
finite cost is submultiplicative. For $q=4$, the stronger conclusion in
Proposition~\ref{prop:full-flip-cup-four} is obtained by tensoring
explicit support-packing certificates. That argument is what
justifies an all-copy obstruction; the finite optimal cost alone
does not supply it.

\subsection{A moment bound for flip and cup coding}

A natural asymmetric extension of the one-copy decomposition keeps
one summand in the algebra of flips and the other in the algebra of
cups. This allows arbitrary coefficient choices, including coding
and threshold rules. The next result excludes all such rules for
$q=4$. Its rank argument is the classical character-moment mechanism
behind the Rao bound for orthogonal arrays
\cite{Rao1947}; see also \cite[Section~3.4 and Theorem~4.1]{Stinson2018}.
The proof below permits arbitrary complex weights and gives the
specific operator-rank consequence needed here.

\begin{proposition}[All-copy bound for two separate operator algebras]
\label{prop:flip-cup-separated-barrier}
Let $q\ge2$, $d=q^2$, $b=q(q-1)/2$, and suppose
\[
 I=M+\Gamma N,\qquad
 M\in\operatorname{span}\{I,F\}^{\otimes k},\qquad
 N\in\operatorname{span}\{I,\Omega\}^{\otimes k}.
\]
Put $C=\operatorname{rank}M+\operatorname{rank}N$.
For every integer $0\le L\le k$,
\[
 C<(d-1)^{k-L}
 \quad\Longrightarrow\quad
 \operatorname{rank}M\ge
 b^k\sum_{j=0}^{\lfloor L/2\rfloor}\binom{k}{j}.
\]
Consequently every sequence of such decompositions, with coefficients
chosen arbitrarily at each $k$, satisfies, for $0<\rho<1/2$,
\[
 \liminf_{k\to\infty}C_k^{1/k}
 \ge\min\left\{(d-1)^{1-2\rho},
       \frac{b}{\rho^\rho(1-\rho)^{1-\rho}}\right\}.
\]
In particular, for $q=4$ and $\rho=1/10$ this lower rate is
\[
 \min\{15^{4/5},60/9^{9/10}\}>8>4^{\sqrt2}.
\]
This concerns the two separate algebras displayed above, rather than
arbitrary pairs in $\mathcal B_q^{\otimes k}$.
\end{proposition}

\begin{proof}
For $S\subseteq[k]$, write $F_S$ for the tensor word with $F$ in
positions $S$ and $I$ elsewhere, and define $\Omega_S$ likewise.
Unique expansion in the flip algebra gives
\[
 M=\sum_S c_SF_S,\qquad N=I-\sum_Sc_S\Omega_S.
\]
On the cup eigenspace in exactly the positions $T\subseteq[k]$,
and its kernel elsewhere, $N$ has eigenvalue
\[
 n(T)=1-\sum_{S\subseteq T}c_Sq^{|S|}.
\]
This block has dimension $(d-1)^{k-|T|}$. The strict cost hypothesis
therefore forces $n(T)=0$ for every $|T|\le L$. Induction on $|T|$
gives
\[
 c_\varnothing=1,\qquad c_S=0\quad(1\le|S|\le L).
\]
On the flip eigenspaces, label signs by $x\in\{-1,1\}^k$ and put
$f(x)=\sum_Sc_S\chi_S(x)$, where
$\chi_S(x)=\prod_{i\in S}x_i$. Every nonzero $f(x)$ contributes
at least $b^k$ to $\operatorname{rank}M$.

Let $t=\lfloor L/2\rfloor$ and index a matrix by subsets of size
at most $t$. Character orthogonality gives
\[
 H_{S,T}=2^{-k}\sum_x f(x)\chi_S(x)\chi_T(x)
        =c_{S\mathbin\triangle T}=\mathbf1_{S=T}.
\]
Thus $H$ is the identity matrix of size
$\sum_{j=0}^{t}\binom{k}{j}$. On the other hand it is a sum of
$|\operatorname{supp}f|$ rank-one matrices,
\[
 H=\sum_{x\in\operatorname{supp}f}
     2^{-k}f(x)\,v(x)v(x)^{\mathsf T},
 \qquad v(x)=(\chi_S(x))_{|S|\le t}.
\]
Its rank is at most that support size, even when the coefficients
are complex. Multiplication by $b^k$ proves the finite bound.

Take $L=\lfloor2\rho k\rfloor$. For each $k$, either the cost
is at least $(d-1)^{k-L}$ or the finite conclusion applies. Taking
$k$th roots and using the binomial-entropy asymptotic gives the
stated minimum. Finally $15^4=50625>32768=8^5$, while
$5^{10}=9765625>6718464=2^{10}3^8$ proves
$60/9^{9/10}>8$. These exact comparisons establish the claimed
strict gap above the square exponent.
\end{proof}

\subsection{The Hodge algebra and a quaternary Fourier equation}

The four-dimensional Hodge operator enlarges the local flip--cup
algebra. The following exact reformulation keeps every coefficient
in the enlarged algebra. Its Fourier matrix is the classical
quaternary MacWilliams matrix
\cite[Theorems~5 and~16]{CRSS1998}; the shadow transform and the
positivity hypotheses used for quantum-code bounds are discussed in
\cite[Theorems~1, 4 and Corollary~7]{Rains1999}.

\begin{proposition}[Hodge algebra as a binary tensor transform]
\label{prop:hodge-fourier}
On $\mathbb C^4\otimes\mathbb C^4$ let
\[
 H_{ij,ab}=\frac12\epsilon_{ijab},\qquad
 \mathcal H=\operatorname{span}_{\mathbb C}\{I,F,\Omega,H\},
\]
where $\epsilon_{1234}=1$ is the alternating symbol. The algebra
$\mathcal H$ is commutative and has four common eigenspaces labelled
$00,01,10,11$, of dimensions $1,3,3,9$. On their scalar eigenvalues
partial transpose acts as $T\otimes T$, where
\[
 T=\frac12\begin{pmatrix}-1&3\\1&1\end{pmatrix}.
\]
Consequently $\mathcal H^{\otimes k}$ identifies with all complex
functions $a$ on $\{0,1\}^{2k}$, with
\[
 \operatorname{rank}a
   =\sum_{x:a(x)\ne0}3^{|x|},\qquad
 \Gamma a=T^{\otimes2k}a.
\]
No sign or permutation-symmetry condition is imposed on $a$.
\end{proposition}

\begin{proof}
Contraction of two alternating symbols gives
$H^2=(I-F)/2$. Also $FH=HF=-H$, $\Omega H=H\Omega=0$,
$\Omega^2=4\Omega$ and $F\Omega=\Omega$. Thus the four pairwise
orthogonal projectors are
\[
 P_{00}=\frac\Omega4,\quad
 P_{01}=\frac{I-F}4+\frac H2,\quad
 P_{10}=\frac{I-F}4-\frac H2,\quad
 P_{11}=\frac{I+F}2-\frac\Omega4.
\]
They sum to $I$. The first and last images are the cup line and
the symmetric traceless subspace. On the six-dimensional exterior
square, $H^2=I$ and $\operatorname{tr}H=0$, so its two eigenspaces
have dimension three each. These observations establish the stated
dimensions and the algebra decomposition.

Partial transpose fixes $I$, exchanges $F,\Omega$, and sends
$H$ to $-H$. Evaluating the transformed projectors gives, in the
specified order,
\[
 \frac14\begin{pmatrix}
 1&-3&-3&9\\-1&-1&3&3\\-1&3&-1&3\\1&1&1&1
 \end{pmatrix}=T\otimes T.
\]
Tensoring these projectors proves both the function-space
identification and the weighted-support rank formula.
\end{proof}

Call a function on $\mathbb F_4^n$ \emph{coordinate-radial} if it
depends only on which coordinates are zero. Here $\mathbb F_4$ is
used as the additive group $\mathbb F_2^2$, with its nondegenerate
alternating form $B$. Define the normalized Fourier transform
\[
 (\mathcal Fg)(y)=2^{-n}\sum_{x\in\mathbb F_4^n}
 (-1)^{\sum_iB(x_i,y_i)}g(x).
\]

\begin{corollary}[An exact support formulation]
\label{cor:hodge-fourier-sparsity}
Put $n=2k$ and $s(x)=(-1)^{\#\{i:x_i=0\}}$.
The equation $A+\Gamma A=I$ in $\mathcal H^{\otimes k}$ is
exactly equivalent to
\[
 g+\mathcal Fg=s
\]
for a coordinate-radial complex function $g$, and
$\operatorname{rank}A=|\operatorname{supp}g|$.
Likewise $M+\Gamma N=I$ is equivalent to
$h+\mathcal Fj=s$, with both functions coordinate-radial and
cost $|\operatorname{supp}h|+|\operatorname{supp}j|$.
\end{corollary}

\begin{proof}
On a single coordinate, a radial function with values $(a,b)$
at zero and at the three nonzero elements transforms to
$((a+3b)/2,(a-b)/2)$. Hence the radial transform is
$K^{\otimes n}$, where
\[
 K=\frac12\begin{pmatrix}1&3\\1&-1\end{pmatrix},\qquad
 D=\begin{pmatrix}-1&0\\0&1\end{pmatrix},\qquad T=-DKD.
\]
Since $n$ is even, $T^{\otimes n}=D^{\otimes n}
K^{\otimes n}D^{\otimes n}$. Multiplication by $D^{\otimes n}$
is multiplication by $s$. Apply this operation to the equation
in Proposition~\ref{prop:hodge-fourier}. A binary word $x$ has
$3^{|x|}$ preimages under the zero/nonzero map, so the weighted
support becomes ordinary support cardinality. The same calculation
proves the asymmetric statement. It also gives $\mathcal Fs=s$.
\end{proof}

This formulation does not import the nonnegativity conditions for
code enumerators. For example, an additive self-dual code $C$ has
$\mathcal F\mathbf1_C=\mathbf1_C$ by character orthogonality,
since $|C|=2^n$. Its indicator therefore gives
$2\mathbf1_C$ on the left, which cannot equal the everywhere
nonzero sign function $s$. Its coordinate-radial average is
nonnegative and likewise cannot equal $s/2$. This excludes that
direct substitution only; linear combinations of code and shadow
data require separate analysis. The support reformulation supplies
no multiplicativity assertion for its minimum cost.
